\clearpage
\renewcommand{\thefigure}{А.\arabic{figure}}
\setcounter{figure}{0}
\section*{ДОДАТОК А}
\begin{center}\bfseries Детальні діаграми класів\end{center}
\addcontentsline{toc}{section}{ДОДАТОК А Детальні діаграми класів}

\fig[\textwidth]{class_core.png}{Діаграма класів пакета \code{core}}
\fig[\textwidth]{class_net_api.png}{Діаграма класів пакетів \code{net}, \code{api},
\code{llm}, \code{search}}

\clearpage
\section*{ДОДАТОК Б}
\begin{center}\bfseries Репозиторій і відтворення результатів\end{center}
\addcontentsline{toc}{section}{ДОДАТОК Б Репозиторій і відтворення результатів}

Вихідний код, історія комітів, UML-джерела, сирі результати вимірювань і ця записка
зберігаються в репозиторії:
\url{https://github.com/d1aboli/course_work_parallel_computing}.

Відтворення всіх результатів (Git Bash, у корені репозиторію, архів \code{D1.7GB.zip} поруч):
\begin{lstlisting}
docker compose up --build -d        # dataset-init, ollama-init, ollama, server
./search --index 12                 # побудова індексу
./search --expand "storm at sea"    # перевірка пошуку
bash bench/run_all.sh               # дослідження, результати в bench/results
python bench/plot.py                # графіки в report/figures
python report/gen_tables.py         # таблиці записки
cd report && xelatex main.tex && xelatex main.tex
\end{lstlisting}

Структура репозиторію:
\begin{lstlisting}
src/main/java/ua/kpi/pocw/
  core/    ThreadPool, Indexer, IndexBuilder, PartitionedIndexBuilder,
           StripedIndexBuilder, TreiberStack, InvertedIndex, BatchTermTable
  net/     HttpServer, HttpRequest, HttpResponse, WebSocket
  api/     Router, SearchService, WsSearchHandler
  search/  Searcher (BM25), Snippets
  llm/     OllamaClient, QueryExpander
  cli/     SearchCli           bench/  Bench
src/test/java/...              власний тест-раннер і 38 тестів
web/                           Web UI
docs/uml/                      PlantUML-діаграми
tools/prepare_dataset.py       підготовка корпусу
bench/                         сценарій дослідження, результати, графіки
report/                        пояснювальна записка (LaTeX)
\end{lstlisting}

\clearpage
\section*{ДОДАТОК В}
\begin{center}\bfseries Промпти, використані під час виконання роботи\end{center}
\addcontentsline{toc}{section}{ДОДАТОК В Промпти, використані під час виконання роботи}

\textit{Промпт 1. Розробка системи}

Мені потрібно зробити курсову з паралельних обчислень: пошукову систему, яка будує
інвертований індекс з текстових файлів у кілька потоків і шукає по ньому. Писати хочу на
Java, без сторонніх бібліотек: пул потоків і структуру індексу для одночасного запису з
багатьох потоків треба написати самому. Система має запускатися однією командою
\code{docker compose up} разом з Ollama і мати REST API та WebSocket за специфікацією з
силабусу, який я додаю.

Поясни, як влаштувати індекс, щоб потоки не заважали один одному, і зроби два варіанти: з
блокуваннями і без них, щоб їх можна було порівняти. Додай простий веб-інтерфейс і консольний
клієнт, тести на відсутність станів гонки й README з інструкцією запуску.

\bigskip
\textit{Промпт 2. Дослідження швидкодії}

Система вже працює на моєму комп'ютері (Ryzen 5 7500F, 6 ядер, 12 потоків). Допоможи
виміряти, як час побудови індексу залежить від кількості потоків від 1 до 16. Кожен варіант
запусти по три рази в перемішаному порядку, щоб результати були чесними. Окремо потрібно
виміряти, скільки запитів за секунду витримує пошук при різній кількості одночасних клієнтів, з
розширенням запиту через LLM і без нього.

Побудуй з результатів графіки й поясни простими словами, де вузьке місце і чому прискорення не
росте далі. Якщо по вимірах видно, що щось можна оптимізувати, зроби це й порівняй «до» і
«після».
